\clearpage\label{answers}
{\sans\bfseries\LARGE\color{ink} Answers}\par
\small
Try every puzzle first. Where an answer shows a calculation, check the steps, not just the result.

\subhead{Chapter 1}
\textbf{Multix on hand spans.} Spans are different sizes for different people. Six of \emph{Multix's} spans is not six of \emph{yours}, so the desk is not ``66~cm'' unless the span is agreed. \textbf{Keep box.} Kept: the number and unit. Set aside: small gaps, overlaps and judgements.

\textbf{1.1} (a) A hides how the book was measured. (b) B tells the most: unit, instrument, edge to edge, nearest mm. (c) A has no unit. 21 what?\\
\textbf{1.2} A 5.3~cm, B 8.8~cm, C 11.6~cm.\\
\textbf{1.3} Dev probably started the pencil at the 1~cm mark and read the other end at 18.4. Subtracting the 1~cm gives 17.4~cm, which fits the others. To check, measure again starting at zero, or use a second ruler.\\
\textbf{1.4} (a) mm ruler. (b) long tape measure. (c) long tape measure.\\
\textbf{1.5} Route 1: $3.5+2.0+4.5+1.5=11.5$~cm. Route 2: $2.5+3.5+2.0+3.0=11.0$~cm. Route 2 is shorter, by 0.5~cm, which is 5~mm.\\
\textbf{1.6} (a) 25~mm, because 25~mm is 2.5~cm. (b) 0.8~m, because it is 80~cm. (c) 3~cm, because it is 30~mm.

\subhead{Chapter 2}
\textbf{Multix on newspapers.} No. ``12\,000'' is a rounded figure, so the real number could be anywhere from about 11\,500 to 12\,499. \textbf{Multix on digits.} More digits is not more exact. The ruler can only show millimetres, so any digit beyond that is made up. \textbf{Keep box.} Kept: about how long each thing is, and the total to the nearest whole. Set aside: the tenths of a centimetre, which add up.

\textbf{2.1} (a) 23 (b) 24 (c) 7 (d) 100 (e) 13 (f) 50 (g) 80 (h) 140.\\
\textbf{2.2} (a) 5.7~cm (b) 4.7~cm (c) 6.3~cm (halves round up).\\
\textbf{2.3} 8.5 $\to$ 8 is wrong. It should be 9 (halves round up).\\
\textbf{2.4} (a) $1.4\times3=4.2$~cm. (b) $1+1+1=3$~cm. (c) 4.2 is closer. Each rounding loses 0.4~cm, three times over, so the pieces lose 1.2~cm together. Round at the end.\\
\textbf{2.5} 8.7 $\to$ 9 (kept: about how long; lost: the 0.3). 46 $\to$ 50 (kept: about how big; lost: the 4). 2.35 $\to$ 2.4 (kept: about how long; lost: the 0.05). Other wordings are fine if they name what was kept and what was lost.\\
\textbf{2.6} The path is 47, 52, 54, 49, 45, because each of these rounds to 50.

\subhead{Chapter 3}
\textbf{Multix on the jar.} 75 is the total of the three scoops, not the jar. First find the average scoop, $75\div3=25$, then multiply by 9 cups: $25\times9=225$. \textbf{Keep boxes.} Beans: kept the typical scoop (the average), set aside which scoop was which and how much they differed. Pacing: kept the typical pace length, set aside the small differences between paces.

\textbf{3.1} Priya 41 off, Sam 29 off, Tomas 6 off. Tomas was closest.\\
\textbf{3.2} (a) $18+20+19+19=76$. (b) $76\div4=19$. (c) $19\times12=228$.\\
\textbf{3.3} 52. It is about double the others. Perhaps two scoops were counted together, or the cup was heaped. Recount that scoop, and write down what you find.\\
\textbf{3.4} (a) A pencil is about 17.5~cm, so 175 has a decimal slip. (b) A door is about 2~m, not 2~mm: a unit slip. (c) A shoe is about 25~cm, not 25~m: a unit slip.\\
\textbf{3.5} The practice field has B $=24$, C $=25$, F $=19$. Total 68, average $68\div3\approx22.67$, times 8 is about 181. The true count is 184, so the estimate is 3 short. Your counts of a square may differ by one or two, and that is fine.

\subhead{Chapter 4}
\textbf{Multix on deleting.} A scientist does not delete readings just because they are awkward. Doing so hides evidence. If a reading is set aside, the reason is written down. \textbf{Keep box.} Kept: the typical time and how much the readings varied. Set aside: the order of the readings and which was which.

\textbf{4.1} Total 65, average 13, spread $15-11=4$.\\
\textbf{4.2} (a) 41. (b) The decimal point was missed; it should be 4.1. (c) $4.1+4.0+4.2+4.1+4.1=20.5$, and $20.5\div5=4.1$.\\
\textbf{4.3} (a) Both averages are 3.0. (b) P: 0.2, Q: $3.8-2.2=1.6$. (c) Class P's readings agree better.\\
\textbf{4.4} No. Her clock reads about 103 seconds every time, so the average is steady but too long. A steady average shows the result can be repeated, not that the clock is right.\\
\textbf{4.5} (a) 25.1 (b) $25.1\div8=3.1375$ (c) 3.1 (d) $3.3-3.0=0.3$.\\
\textbf{Catch the ruler (practice numbers).} Total 130, average 13.0~cm, spread 6~cm, about one sixth of a second.

\subhead{Chapter 5}
\textbf{Multix on the scale.} No. The apple is $152-3=149$~g. Zero is not just a feeling. \textbf{Keep box.} Kept: the corrected length, with the old reading still visible in the wrap. Set aside: the misleading raw reading as an answer.

\textbf{5.1} A: start 2.0, end 9.4, length 7.4~cm. B: start 1.0, end 12.3, length 11.3~cm. C: start 5.5, end 11.5, length 6.0~cm.\\
\textbf{5.2} (a) $205-5=200$~g. (b) $22.0-1.5=20.5$~cm.\\
\textbf{5.3} C. Reading from the side makes the marks seem to move.\\
\textbf{5.4} (a) $400-150=250$~mL. (b) $35-18=17$~$^\circ$C. (c) One reading cannot show whether the jug is steady. It shows only that this jug reads 25~mL above the better jug this once. Read both again and compare before deciding which one is off.\\
\textbf{The jug and thermometer figure.} 275~mL and 30~$^\circ$C.

\subhead{Chapter 6}
\textbf{6.1} (a) O. (b) C. (c) E. (d) K. The ten readings did not test sentence (d). It is a claim about why, and it needs its own test.\\
\textbf{6.2} Far too many digits for a ruler marked in millimetres. Write 17.4~cm.\\
\textbf{6.3} (a) 123.6 (b) 24.72 (c) 24.7 (d) $24.9-24.6=0.3$. (e) For example: ``The string is about 24.7~cm long; five readings with a mm ruler ranged from 24.6 to 24.9~cm.''\\
\textbf{6.4} (a) measure many times, average. (b) subtract the zero error. (c) look straight on. (d) round to what the tool can show.\\
\textbf{Slop check.} (a) You cannot tell: there is no title and no axis labels. (b) Many fonts and colours, stars, rotated filler text and ``lol''. (c) Axis labels, units, and a title saying what each bar means. (d) The stars, the filler text and the jokes.\\
\textbf{Redraw.} A good version has one title, labelled axes with units, one or two calm colours, and nothing else.

\textbf{Whole-book table.} Chapter 2: simplified a length by rounding; kept ``about how long''; set aside tenths. Chapter 3: simplified a jar by scooping; kept the typical scoop; set aside differences between scoops. Chapter 4: simplified ten readings to an average; kept average and spread; set aside order. Chapter 5: corrected for a start that is not at zero; kept the true length; set aside the misleading raw reading as the answer. Chapter 6: simplified an experiment into five lines; kept the answer with its note; set aside the detail of every step.

\clearpage\label{notes}
{\sans\bfseries\LARGE\color{ink} Words}\par
\small
\begin{multicols}{2}
\word{average}{the total of the readings divided by how many there are. Also called the mean.}
\word{count}{a number from counting, such as beans in a scoop.}
\word{dot plot}{a picture with one dot for each reading.}
\word{estimate}{a number about more than was counted, worked out by reasoning.}
\word{instrument}{a tool for measuring, such as a ruler, scale or stopwatch.}
\word{middle value}{the reading in the middle when the readings are put in order. Also called the median.}
\word{parallax}{an apparent change in a reading when you look from the side.}
\word{round}{to replace a number with a nearby, simpler one.}
\word{spread}{the biggest reading minus the smallest.}
\word{unit}{an agreed size to measure with, such as the centimetre.}
\word{zero error}{when an instrument does not start exactly at zero.}
\end{multicols}

\vspace{0.4em}
{\sans\bfseries\large\color{sky!70!black} Notes for grown-ups}\par
This book teaches that a measurement is a careful judgement, not a perfect fact. Each chapter has one activity at a desk or table, one worked example with made-up practice numbers (labelled as such), one trust box, one ``what did we keep?'' prompt, and a page of puzzles.

\textbf{Materials and safety.} A ruler marked in millimetres, a second ruler, string, a timer showing seconds, a small cup, a jar, a cup of dried beans or rice, scrap paper. Dried beans are a choking risk for little children and pets. Falling-paper drops are done on the floor, never on furniture. A kitchen scale is optional and needs an adult.

\textbf{Check boxes.} The trust boxes say what a check shows and what it does not show. Repeating a measurement shows that it can be repeated, not that it is right. Two rulers agreeing makes a shared fault unlikely, but does not rule it out.

\textbf{The practice numbers} in the worked examples (beans, falls, pencils) were made up so that the arithmetic is easy to check. They are not results of real experiments, and your own results will differ. That is the point of the book.

\subhead{Sources}
\begin{itemize}[leftmargin=1.3em,itemsep=0.2em]
\item Bureau International des Poids et Mesures. \emph{The International System of Units (SI)}, 9th edition, 2019. For the definition of the metre in terms of the speed of light.
\item J. R. Taylor. \emph{An Introduction to Error Analysis}, 2nd edition. University Science Books, 1997. For repeated measurement, spread, and the difference between random and systematic error.
\item B. N. Taylor and C. E. Kuyatt. \emph{Guidelines for Evaluating and Expressing the Uncertainty of NIST Measurement Results}, NIST Technical Note 1297, 1994. For how to report a result with a stated uncertainty.
\end{itemize}

\clearpage
\thispagestyle{blank}
\begin{tikzpicture}[remember picture,overlay]
  \fill[sun!35] ($(current page.south west)+(0.4in,0.4in)$) rectangle ($(current page.north east)-(0.4in,0.4in)$);
  \node[anchor=north,align=center,font=\sffamily\bfseries,text=ink,text width=5.4in] at ($(current page.north)+(0,-1.4in)$) {\Large Next in Science\\[6pt]\fontsize{28}{30}\selectfont A Fair Test};
  \node[align=center,font=\sffamily,text width=4.6in,text=ink] at (current page.center) {Paper planes fly different distances, and the reason is easy to get wrong. Book 2 shows how to change one thing at a time, keep everything else the same, and write it down so a stranger could repeat your test.};
  \node at ($(current page.south)+(0,2.6in)$) {\scalebox{3}{\multixbot[ok]}};
  \node[font=\sffamily\bfseries,text=ink] at ($(current page.south)+(0,0.9in)$) {Ages 9 to 11 \quad$\star$\quad Science Book 1};
\end{tikzpicture}
